\documentclass[11pt]{amsart}
\usepackage[margin=1in]{geometry}
\usepackage{amsmath,amssymb,amsthm}
\usepackage{hyperref}

\newtheorem{theorem}{Theorem}
\newtheorem{lemma}{Lemma}

\newcommand{\C}{\mathbb{C}}
\newcommand{\N}{\mathbb{N}}
\newcommand{\one}{\mathbf{1}}

\title{A Dyadic Counterexample to an Ideal-Closure Question of Glueck}
\author{fa\_banach\_001, agent\_lane\_10}
\date{2026-06-18}

\begin{document}
\maketitle

\section*{Source and Target}

The source paper is Jochen Glueck, \emph{Growth rates and the peripheral
spectrum of positive operators}, arXiv:1512.07483.  Lemma 4.4 of that paper
considers a complex Banach lattice \(E\) and a sequence \((S_n)\) of positive
operators on \(E\). It defines
\[
 I=\{x\in E:\exists x_n\to x \hbox{ such that } S_n|x_n|\to 0\}.
\]
In Remarks 4.5(a)--(c), the paper also considers
\[
 J=\{x\in E:S_n|x|\to 0\}.
\]
It observes that \(J\subset \overline J\subset I\), that equality holds when
\((S_n)\) is norm bounded, and asks whether \(I=\overline J\) holds in general
when \((S_n)\) is not norm bounded.

\section*{Idea of the Proof}

The relaxed ideal \(I\) allows the approximating vector to depend on \(n\).
Thus one can defeat the \(n\)-th operator by modifying the function on only the
small set inspected by that operator. The fixed ideal \(J\), however, asks one
function to defeat all sufficiently fine tests at once.

On \(L^1[0,1]\), let the tests be normalized averages over all dyadic
intervals, enumerated level by level. The constant function \(1\) belongs to
\(I\), because for each dyadic interval \(D\) we can delete \(D\), changing
\(1\) in \(L^1\)-norm by \(|D|\to 0\), and thereby making the corresponding
average equal to zero. But a function in \(J\) must have averages of its
absolute value tending uniformly to zero over every dyadic interval at high
levels. Dyadic differentiation then forces the function to be zero almost
everywhere.

\section*{Counterexample}

\begin{theorem}
There are a complex Banach lattice \(E\) and a sequence \((S_n)\) of positive
operators on \(E\), not bounded in operator norm, for which the ideals \(I\)
and \(J\) from Lemma 4.4 and Remarks 4.5 of Glueck's paper satisfy
\[
   \overline J \subsetneq I.
\]
Indeed, one may take \(E=L^1([0,1];\C)\), in which case \(J=\{0\}\) but
\(\one\in I\).
\end{theorem}

\begin{proof}
Let \(E=L^1([0,1];\C)\) with its usual complex Banach lattice structure.  For
each \(k\in\N\) and \(1\le j\le 2^k\), put
\[
  D_{k,j}=\big[(j-1)2^{-k},j2^{-k}\big),
\]
with the harmless convention that the last interval includes the endpoint
\(1\). Enumerate all pairs \((k,j)\) level by level. If \(n\) corresponds to
\((k,j)\), define \(S_n=S_{k,j}\in\mathcal L(E)\) by
\[
  S_{k,j}f
  =2^k\left(\int_{D_{k,j}} f(t)\,dt\right)\one .
\]
This is a positive rank-one operator. Its norm is \(2^k\), so the sequence is
not norm bounded. Moreover, for \(f\in E\),
\[
  \|S_{k,j}|f|\|_1
  =2^k\int_{D_{k,j}} |f(t)|\,dt,
\]
the average of \(|f|\) over the dyadic interval \(D_{k,j}\).

We first show that \(J=\{0\}\). Suppose \(f\in J\). Since the enumeration is by
dyadic levels, the condition \(\|S_n|f|\|_1\to0\) implies
\[
  \max_{1\le j\le 2^k} 2^k\int_{D_{k,j}} |f(t)|\,dt \longrightarrow 0
  \quad (k\to\infty).
\]
For almost every \(t\in[0,1]\), the dyadic intervals \(D_k(t)\) of level \(k\)
which contain \(t\) differentiate \(L^1\); hence
\[
  2^k\int_{D_k(t)} |f(s)|\,ds \longrightarrow |f(t)|
\]
for almost every \(t\). The preceding uniform convergence of all level-\(k\)
dyadic averages to \(0\) therefore gives \(|f(t)|=0\) for almost every \(t\).
Thus \(f=0\), and \(J=\{0\}\). In particular \(\overline J=\{0\}\).

It remains to show that \(I\) is larger. Let \(n\) correspond to \(D_{k,j}\)
and set
\[
  f_n=\one_{[0,1]\setminus D_{k,j}} .
\]
As \(n\to\infty\), the dyadic level \(k\to\infty\), and hence
\[
  \|f_n-\one\|_1=|D_{k,j}|=2^{-k}\longrightarrow 0.
\]
On the other hand, \(|f_n|=0\) on \(D_{k,j}\), so
\[
  S_n|f_n|=0
\]
for every \(n\). Hence \(\one\in I\). Since \(\one\notin\overline J=\{0\}\),
we have \(\overline J\subsetneq I\).
\end{proof}

\section*{Verification Notes}

The construction uses only normalized dyadic interval averages on
\(L^1[0,1]\). The only external analytic input is the standard dyadic
Lebesgue differentiation theorem. The endpoint convention for dyadic intervals
does not matter because endpoints have measure zero.

The packet includes a small finite-level sanity script which checks the
identity
\[
  \|S_{k,j}|f_{k,j}|\|_1=0,\qquad \|f_{k,j}-\one\|_1=2^{-k},
\]
for the simple functions used in the proof.

\section*{Novelty Check and Limitations}

The run indexes were searched for arXiv:1512.07483, the source title, Lemma
4.4, the phrase \(I\) coincides with \(\overline J\), and nearby
Banach-lattice ideal terminology. No prior packet or attempt was found.

External web searches were also performed for exact and close phrases from
Remarks 4.5(c), including the author's ``does not know whether'' sentence. No
explicit later answer was found in this bounded pass.

This packet does not address the global cyclicity problem for the peripheral
spectrum of positive operators. It answers only the auxiliary ideal-closure
question in Remarks 4.5(c), by counterexample.

\section*{References}

\begin{thebibliography}{9}
\bibitem{Gluck1512}
Jochen Glueck,
\emph{Growth rates and the peripheral spectrum of positive operators},
arXiv:1512.07483.

\bibitem{DyadicDiff}
The dyadic differentiation theorem for \(L^1[0,1]\), equivalently the
Lebesgue differentiation theorem restricted to the dyadic filtration.
\end{thebibliography}

\section*{Human Review Recommendation}

Candidate counterexample, likely valid. Review should focus on whether
Remarks 4.5 use precisely the norm convergence \(\|S_n|x|\|\to0\), and on the
dyadic differentiation step proving \(J=\{0\}\).

\end{document}
